\chapter{径向势与氢原子}
\label{chap:radial-hydrogen}

\section{径向势}
\label{sec:radial-potentials}

若 $V$ 是 $\mathbb R^3$ 上任意径向函数,令
\[
 \widehat H=-\frac{\hbar^2}{2m}\Delta+V
\]
为作用于 $L^2(\mathbb R^3)$ 的相应 Hamilton 算子.我们将寻找定态 Schrödinger 方程 $\widehat H\psi=E\psi$ 的解,解的形式为
\[
 \psi(x)=p(x)f(|x|),
\]
其中 $f$ 是 $(0,\infty)$ 上的光滑函数,$p$ 是 $\mathbb R^3$ 上次数为 $l$ 的齐次调和多项式.

\begin{proposition}[径向因子与 Laplace 算子]\label{prop:radial-laplacian}
设 $p$ 是 $\mathbb R^3$ 上次数为 $l$ 的齐次调和多项式,$f$ 是 $(0,\infty)$ 上的光滑函数.令 $\psi$ 为 $\mathbb R^3\setminus\{0\}$ 上的函数
\[
 \psi(x)=p(x)f(|x|).
\label{eq:radial-laplacian-function}
\]
则在 $\mathbb R^3\setminus\{0\}$ 上有
\[
 \Delta\psi(x)=p(x)\left(\frac{d^2f}{dr^2}+\frac{2(l+1)}r\frac{df}{dr}\right).
\label{eq:radial-laplacian}
\]
\end{proposition}

\begin{proof}
先考虑 $l=0$ 的情形,此时 $p$ 是常数,我们取 $p=1$,而 $\psi$ 只是径向函数 $f(|x|)$.于是
\[
\begin{aligned}
 \frac{\partial}{\partial x_j}f(|x|)
 &=\frac{df}{dr}\frac{d}{dx_j}(x_1^2+x_2^2+x_3^2)^{1/2}\\
 &=\frac{df}{dr}\frac{x_j}{|x|}.
\end{aligned}
\]
因此
\[
\begin{aligned}
 \sum_{j=1}^3\frac{\partial^2}{\partial x_j^2}f(|x|)
 &=\sum_{j=1}^3\left(\frac{d^2f}{dr^2}\frac{x_j^2}{|x|^2}
 +\frac{df}{dr}\left(\frac1{|x|}-\frac{x_j^2}{|x|^3}\right)\right)\\
 &=\frac{d^2f}{dr^2}+\frac2r\frac{df}{dr}.
\end{aligned}
\]

一般情形下,Laplace 算子的乘积法则给出
\[
 \Delta\psi=(\Delta p)f(|x|)+2\nabla p\cdot\nabla f(|x|)+p\Delta f(|x|).
\]
由假设 $\Delta p=0$.此外,由于 $f(|x|)$ 是径向函数,其梯度指向径向方向,所以只需要考虑 $\nabla p$ 的径向分量.而在穿过原点的每条射线上,$p$ 的行为都像某个常数乘以 $r^l$,因此 $p$ 对 $r$ 的导数为 $(l/r)p$.于是
\[
 \Delta\psi=\frac{2l}{r}p\frac{df}{dr}+p\frac{d^2f}{dr^2}+p\frac2r\frac{df}{dr},
\]
这就化为所要证明的式~\eqref{eq:radial-laplacian}.
\end{proof}

\def\ChapterEighteenLateCalculations{
\begin{lemma}[Runge--Lenz 动量项的交换关系]\label{lem:runge-lenz-momentum-cross-commutator}
对所有 $j,m$,有
\[
 [(P\times\widehat J)_j,(P\times\widehat J)_m]
 =-i\hbar(P\cdot P)\varepsilon_{jml}\widehat J_l.
\label{eq:runge-lenz-momentum-cross-commutator}
\]
\end{lemma}

\begin{proof}
计算 $[P_k\widehat J_l,P_n\widehat J_o]$ 时,反复使用交换子的乘积法则(命题~\ref{prop:commutator-properties} 的第 3 点).会得到四项,其中一项(含 $[P_k,P_n]$ 的项)为零.对其余各项使用式~\eqref{eq:vector-operator-epsilon-commutator},得到
\[
\begin{aligned}
 \frac1{i\hbar}[\varepsilon_{jkl}P_k\widehat J_l,
 \varepsilon_{mno}P_n\widehat J_o]
 =\varepsilon_{jkl}\varepsilon_{mno}\bigl(&P_k[\widehat J_l,P_n]\widehat J_o\\
 &+P_nP_k[\widehat J_l,\widehat J_o]
 +P_n[P_k,\widehat J_o]\widehat J_l\bigr).
\end{aligned}
\label{eq:runge-lenz-pj-commutator-expanded}
\]
先计算式~\eqref{eq:runge-lenz-pj-commutator-expanded} 右端的第一项.利用引理~\ref{lem:epsilon-identities} 和 $P$ 是向量算子这一事实,
\[
\begin{aligned}
 \varepsilon_{jkl}\varepsilon_{mno}P_k[\widehat J_l,P_n]\widehat J_o
 &=\varepsilon_{jkl}(\delta_{op}\delta_{ml}-\delta_{ol}\delta_{mp})P_kP_p\widehat J_o\\
 &=\varepsilon_{jkm}P_kP_p\widehat J_p-\varepsilon_{jko}P_kP_m\widehat J_o\\
 &=\varepsilon_{jkm}P_k(P\cdot\widehat J)-P_m(P\times\widehat J)_j.
\end{aligned}
\]
类似地计算第二项和第三项,得到
\[
\begin{aligned}
 \frac1{i\hbar}[\varepsilon_{jkl}P_k\widehat J_l,
 \varepsilon_{mno}P_n\widehat J_o]
 ={}&\varepsilon_{jkm}P_k(P\cdot\widehat J)-P_m(P\times\widehat J)_j\\
 &+(P\times P)_j\widehat J_m-\varepsilon_{jkm}P_k(P\cdot\widehat J)\\
 &+P_m(P\times\widehat J)_j-(P\cdot P)\varepsilon_{jml}\widehat J_l.
\end{aligned}
\]
其中涉及 $P\cdot\widehat J$ 或 $P\times P$ 的三项为零,另有两项相消,剩下
\[
 \frac1{i\hbar}[\varepsilon_{jkl}P_k\widehat J_l,
 \varepsilon_{mno}P_n\widehat J_o]
 =-(P\cdot P)\varepsilon_{jml}\widehat J_l,
\]
即得结论.
\end{proof}

现在继续证明 Runge--Lenz 向量的性质.

\begin{proof}[命题~\ref{prop:quantum-runge-lenz-algebraic} 的证明]
由引理~\ref{lem:runge-lenz-orthogonality} 中的第一组正交关系,立即得到
\[
 \widehat J\cdot\widehat A=\widehat A\cdot\widehat J=0.
\]
另一方面,使用式~\eqref{eq:quantum-runge-lenz-alternative} 表达 $\widehat A$,展开 $\widehat A\cdot\widehat A$ 并稍作整理,得到
\[
\begin{aligned}
 \widehat A\cdot\widehat A
 ={}&1+\frac1{\mu^2Q^4}(P\cdot P)(\widehat J\cdot\widehat J+\hbar^2)\\
 &-\frac1{\mu Q^2}\left(2\widehat J\cdot\widehat J
 +i\hbar\left[P\cdot\frac{X}{|X|}-\frac{X}{|X|}\cdot P\right]\right).
\end{aligned}
\]
现在,由正则对易关系,
\[
\begin{aligned}
 \frac{X}{|X|}\cdot P-P\cdot\frac{X}{|X|}
 &=i\hbar\left(\frac{\delta_{kk}}{|X|}-\frac{X_kX_k}{|X|^3}\right)\\
 &=\frac{2i\hbar}{|X|}.
\end{aligned}
\]
因此
\[
 \widehat A\cdot\widehat A
 =1+(\widehat J\cdot\widehat J+\hbar^2)
 \left(\frac{P\cdot P}{\mu^2Q^4}-\frac{2}{\mu Q^2|X|}\right).
\]
注意
\[
 \frac{P\cdot P}{\mu^2Q^4}-\frac{2}{\mu Q^2|X|}
 =\frac{2\widehat H}{\mu Q^4},
\]
故得到式~\eqref{eq:quantum-runge-lenz-square}.
\end{proof}

\begin{proof}[引理~\ref{lem:quantum-runge-lenz-commutator} 的第二部分]
使用式~\eqref{eq:quantum-runge-lenz-alternative} 表达 $\widehat A$.计算 $\widehat A_j$ 与 $\widehat A_m$ 的交换子时,会出现几类不同的项,我们逐类计算.显然,$X_j/|X|$ 与 $X_m/|X|$ 的交换子为零.$P\times\widehat J$ 两项之间的交换子已经由引理~\ref{lem:runge-lenz-momentum-cross-commutator} 计算.

另一方面,计算 $[P_k\widehat J_l,X_m/|X|]$ 时同样得到四项,其中一项为零,即含有 $[\widehat J_l,1/|X|]$ 的项,因为 $1/|X|$ 在旋转下不变.于是
\[
\begin{aligned}
 \frac1{i\hbar}\left[\varepsilon_{jkl}P_k\widehat J_l,\frac{X_m}{|X|}\right]
 ={}&-\varepsilon_{jml}\widehat J_l\frac1{|X|}
 +\varepsilon_{jkl}\varepsilon_{lmn}P_kX_n\frac{X_m}{|X|^3}\\
 &+\varepsilon_{jkl}X_m\frac{X_k}{|X|^3}\varepsilon_{lno}X_nP_o.
\end{aligned}
\]
应用引理~\ref{lem:epsilon-identities} 并进行与前面类似的计算,得到
\[
\begin{aligned}
 \frac1{i\hbar}\left[\varepsilon_{jkl}P_k\widehat J_l,\frac{X_m}{|X|}\right]
 ={}&-\varepsilon_{jml}\widehat J_l\frac1{|X|}
 +\delta_{jm}\frac{P\cdot X}{|X|}\\
 &+X_mX_j\frac{X\cdot P}{|X|^5}
 -P_m\frac{X_j}{|X|}+\frac{X_m}{|X|}P_j.
\end{aligned}
\label{eq:runge-lenz-cross-position-commutator}
\]

在交换子 $[\alpha_j+\beta_j,\alpha_m+\beta_m]$ 中,涉及一个 $\alpha$ 和一个 $\beta$ 的项为 $[\alpha_j,\beta_m]+[\beta_j,\alpha_m]=[\alpha_j,\beta_m]-[\alpha_m,\beta_j]$.这个量关于 $j,m$ 反对称,即交换 $j,m$ 会改变符号.因此,式~\eqref{eq:runge-lenz-cross-position-commutator} 中关于 $j,m$ 对称的项,在计算 $[\widehat A_j,\widehat A_m]$ 的完整表达式时都会消失,所以第二、第三项可以忽略.对于最后一项,把 $P_m$ 移过 $X_j/|X|$,得到
\[
 \frac{X_j}{|X|}P_m+\frac{X_m}{|X|}P_j
 =\frac{X_j}{|X|}P_m+\frac{X_m}{|X|}P_j
 -i\hbar\left(\frac{\delta_{jm}}{|X|}-\frac{X_jX_m}{|X|^3}\right),
\label{eq:runge-lenz-symmetric-cross-term}
\]
这也是关于 $j,m$ 对称的.因此,式~\eqref{eq:runge-lenz-cross-position-commutator} 中只有第一项会对 $[\widehat A_j,\widehat A_m]$ 的计算作出贡献;这一项关于 $j,m$ 反对称,在完整交换子中会出现两次.

现在直接计算 $[\varepsilon_{jkl}P_k\widehat J_l,P_m]$ 与 $[P_j,X_m/|X|]$,可以验证它们关于 $j,m$ 对称(习题 8),因此不会贡献于 $[\widehat A_j,\widehat A_m]$.最后得到
\[
\begin{aligned}
 \frac1{i\hbar}[\widehat A_j,\widehat A_m]
 &= -\frac1{\mu^2Q^4}\varepsilon_{jml}(P\cdot P)\widehat J_l
 +\frac1{\mu^2Q^4}\varepsilon_{jml}\widehat J_l\frac{Q^2}{|X|}\\
 &= -\frac2{\mu Q^4}\varepsilon_{jml}\widehat J_l
 \left(\frac{P\cdot P}{2\mu}-\frac{Q^2}{|X|}\right)\\
 &= -\frac2{\mu Q^4}\varepsilon_{jml}\widehat J_l\widehat H,
\end{aligned}
\]
这正是引理中所声称的关系.
\end{proof}

\begin{proof}[定理~\ref{thm:quantum-so4-symmetry} 的证明]
由于 Hamilton 算子 $\widehat H$ 在旋转下不变,$\widehat H$ 与角动量的每个分量对易;我们也已经证明 $\widehat H$ 与 Runge--Lenz 向量的每个分量对易.因此 $\widehat I,\widehat K$ 与 Hamilton 算子对易.

由于 $\widehat A_k$ 与 $\widehat H$ 对易,它也与 $\widehat H$ 的任意函数对易.由引理~\ref{lem:quantum-runge-lenz-commutator},
\[
\begin{aligned}
 \frac1{i\hbar}[\widehat B_j,\widehat B_k]
 &=\frac{\mu Q^4}{2|\widehat H|}\frac1{i\hbar}[\widehat A_j,\widehat A_k]\\
 &=-\frac{\mu Q^4}{2|\widehat H|}\frac2{\mu Q^4}
 \varepsilon_{jkl}\widehat J_l\widehat H.
\end{aligned}
\]
在负能量子空间 $V^-$ 上,$\widehat H/|\widehat H|=-I$,所以这个表达式化为 $\varepsilon_{jkl}\widehat J_l$.(正能量子空间上的结果会与 $V^-$ 上的结果相差一个关键的负号.)

另一方面,由于 $\widehat B$ 与 $\widehat J$ 都是向量算子,命题~\ref{prop:vector-operator-commutators} 给出
\[
 \frac1{i\hbar}[\widehat B_j,\widehat J_k]=\varepsilon_{jkl}\widehat B_l,
 \qquad
 \frac1{i\hbar}[\widehat J_j,\widehat J_k]=\varepsilon_{jkl}\widehat J_l.
\]
利用 $\widehat B_j$ 与 $\widehat J_j$ 的这些交换关系,直接计算即可验证 $\widehat I$ 与 $\widehat K$ 的各分量满足定理中的交换关系.最后,式~\eqref{eq:quantum-runge-lenz-so4-casimir} 由命题~\ref{prop:quantum-runge-lenz-algebraic} 中的平方关系和 $\widehat I,\widehat K$ 的定义直接得到.
\end{proof}
}

\def\ChapterEighteenExercises{
\section{练习}
\label{sec:chapter18-exercises}

\begin{exercise}[两粒子氢样 Hamilton 算子]\label{ex:hydrogen-two-particle-reduction}
考虑 $\mathbb R^3$ 中两个粒子通过 $1/r$ 势相互作用时的量子 Hamilton 算子
\[
 \widehat H=-\frac{\hbar^2}{2m_1}\Delta_1-\frac{\hbar^2}{2m_2}\Delta_2-\frac{Q^2}{|x_1-x_2|}.
\]
这里,如第~\ref{sec:quantum-particle-rn} 节,$\Delta_1,\Delta_2$ 分别是关于变量 $x_1,x_2$ 的 Laplace 算子.像第~\ref{sec:particle-systems} 节那样,引入质心
\[
 c=\frac{m_1x_1+m_2x_2}{m_1+m_2}
\]
和相对位置 $y=x_1-x_2$.证明在这些新变量中,$\widehat H$ 可以写为
\[
 -\frac{\hbar^2}{2(m_1+m_2)}\Delta_c
 -\frac{\hbar^2}{2\mu}\Delta_y-\frac{Q^2}{|y|},
\]
其中 $\mu=m_1m_2/(m_1+m_2)$ 是约化质量.

\begin{note}
在新变量中,$\widehat H$ 是两个项之和:一项只含变量 $c$,另一项只含变量 $y$.只含 $c$ 的项是质量为 $m_1+m_2$ 的自由粒子的 Hamilton 算子;只含 $y$ 的项是质量为 $\mu$、处在 $1/r$ 势中的粒子的 Hamilton 算子.
\end{note}
\end{exercise}

\begin{exercise}[负能量经典 Kepler 区域的有限体积]\label{ex:negative-energy-kepler-finite-volume}
令 $H(x,p)=|p|^2/(2\mu)-Q^2/|x|$ 为 $\mathbb R^3$ 中经典 Kepler 问题的 Hamilton 函数.证明对每个 $\varepsilon>0$,$\mathbb R^6$ 中由
\[
 \{(x,p)\mid H(x,p)<-\varepsilon\}
\]
给出的区域具有有限体积.
\end{exercise}

\begin{exercise}[四元数代数]\label{ex:quaternion-algebra}
令 $\mathbb H$ 表示 $M_2(\mathbb C)$ 中以下四个元素的实线性张成空间:
\[
 1=\begin{pmatrix}1&0\\0&1\end{pmatrix},\quad
 i=\begin{pmatrix}i&0\\0&-i\end{pmatrix},\quad
 j=\begin{pmatrix}0&1\\-1&0\end{pmatrix},\quad
 k=\begin{pmatrix}0&i\\i&0\end{pmatrix}.
\]
\begin{enumerate}
\item 证明在矩阵乘法下,$\mathbb H$ 是一个实结合代数,并满足
\[
 i^2=j^2=k^2=-1,\qquad ij=-ji=k,\qquad jk=-kj=i,\qquad ki=-ik=j.
\]
这个代数 $\mathbb H$ 是四元数代数的一种具体实现.
\item 证明 $\mathbb H$ 的每个非零元素都有乘法逆元.
\end{enumerate}
提示:仿照每个非零复数都有乘法逆元的论证.
\end{exercise}

\begin{exercise}[$\mathfrak{so}(4)$ 与两个 $\mathfrak{so}(3)$ 的同构]\label{ex:so4-so3-direct-sum}
令 $\mathbb H$ 表示习题 3 中定义的四元数代数.本题显式建立 Lie 代数 $\mathfrak{so}(4)$ 与 $\mathfrak{so}(3)\oplus\mathfrak{so}(3)$ 之间的同构(比较定义~\ref{def:representation-sum-tensor-dual}).
\begin{enumerate}
\item 令 $V$ 为由 $i,j,k$ 张成的子空间.证明关于括号 $[\alpha,\beta]=\alpha\beta-\beta\alpha$,$V$ 构成 Lie 代数,并且与 $\mathfrak{so}(3)$ 同构.
\item 令 $\operatorname{End}(\mathbb H)$ 表示 $\mathbb H$ 到自身的实线性映射代数.对 $\alpha\in V$,令 $L_\alpha(\beta)=\alpha\beta$ 为“左乘 $\alpha$”映射,令 $R_\alpha(\beta)=\beta\alpha$ 为“右乘 $\alpha$”映射.证明映射 $\alpha\mapsto L_\alpha$ 与 $\alpha\mapsto-R_\alpha$ 都是从 $V$ 到 $\operatorname{End}(\mathbb H)$ 的 Lie 代数同态.
\item 在 $\mathbb H$ 上取使 $\{1,i,j,k\}$ 成为标准正交基的内积.对 $\alpha\in V$,证明
\[
 \langle L_\alpha\beta,\gamma\rangle=-\langle\beta,L_\alpha\gamma\rangle,
 \qquad
 \langle R_\alpha\beta,\gamma\rangle=-\langle\beta,R_\alpha\gamma\rangle.
\]
也就是说,$L_\alpha,R_\alpha$ 属于 $\mathfrak{so}(4)$;这里把 $\mathfrak{so}(4)$ 识别为关于该内积斜对称的 $\operatorname{End}(\mathbb H)$ 元素空间.
\item 证明映射 $(\alpha,\beta)\mapsto L_\alpha-R_\beta$ 是从 $\mathfrak{so}(3)\oplus\mathfrak{so}(3)$ 到 $\mathfrak{so}(4)$ 的 Lie 代数同构.
\item 令 $D$ 表示 $\mathfrak{so}(3)\oplus\mathfrak{so}(3)$ 的对角子代数,即形如 $(X,X)$ 的元素全体.证明第 4 部分同构下 $D$ 的像,是相对于第 3 部分基具有形式
\[
 Y=\begin{pmatrix}0&0\\0&Z\end{pmatrix},
 \qquad Z\in\mathfrak{so}(3)
\]
的 $\mathfrak{so}(4)\subset\operatorname{End}(\mathbb H)$ 元素全体.
\end{enumerate}
\end{exercise}

\begin{exercise}[\(\mathfrak{so}(4)\) 中的两个子代数]\label{ex:so4-subalgebras}
在习题 4 的同构
\[
 \mathfrak{so}(4)\cong\mathfrak{so}(3)\oplus\mathfrak{so}(3)
\]
下,显式描述对应于两个 $\mathfrak{so}(3)$ 的两个子代数.
\end{exercise}

\begin{exercise}[验证 Levi--Civita 恒等式]\label{ex:verify-epsilon-identities}
验证引理~\ref{lem:epsilon-identities}.

提示:先证明 $\varepsilon_{jkl}\varepsilon_{jmn}$ 只有在 $(k,l)=(m,n)$ 或 $(k,l)=(n,m)$ 时才可能非零.
\end{exercise}

\begin{exercise}[向量算子叉积的变换性质]\label{ex:vector-operator-cross-product}
本题使用第~\ref{subsec:runge-lenz-notation} 节的求和约定.
\begin{enumerate}
\item 证明对任意 $3\times3$ 矩阵 $M$ 以及 $j,k,l\in\{1,2,3\}$,有
\[
 \varepsilon_{mno}M_{jm}M_{kn}M_{lo}=\varepsilon_{jkl}\det M.
\]
\item 证明若 $C$ 是向量算子,则对所有 $R\in SO(3)$,
\[
 \Pi(R)C_k\Pi(R)^{-1}=R_{lk}C_l.
\]
\item 证明两个向量算子的叉积仍是向量算子.
\end{enumerate}
提示:把向量算子的定义写成等价形式
\[
 v\cdot C=\Pi(R)((R^{-1}v)\cdot C)\Pi(R)^{-1}.
\]
\end{exercise}

\begin{exercise}[对称交换子计算]\label{ex:runge-lenz-symmetric-commutators}
计算 $[\varepsilon_{jkl}P_k\widehat J_l,P_m]$ 与 $[P_j,X_m/|X|]$,并证明这两个量关于 $j,m$ 对称,即交换 $j,m$ 后取值不变.
\end{exercise}

\begin{exercise}[径向方程的两个幂级数解]\label{ex:hydrogen-radial-power-series}
证明式~\eqref{eq:hydrogen-laguerre-equation} 对 $g(\rho)$ 有两个幂级数解:一个以 $\rho^{-(2l+1)}$ 开头,另一个以 $\rho^0$ 开头.
\end{exercise}
}

\def\ChapterEighteenClassical{
\section{量子 Kepler 问题中的 Runge--Lenz 向量}
\label{sec:quantum-runge-lenz}

在第~\ref{sec:kepler-problem} 节中,我们利用作为守恒量的 Runge--Lenz 向量,几乎完全解决了经典 Kepler 问题.Runge--Lenz 向量的量子版本与 Hamilton 算子对易,并且可以说明量子 Kepler 问题(通常把它看作氢原子问题)的许多特殊性质.特别地,Runge--Lenz 向量将帮助我们解释:
\begin{enumerate}
\item 氢原子的负能量为何具有 $-R/n^2$ 的简单形式;
\item 对给定的 $n$,式~\eqref{eq:hydrogen-bound-state} 中各状态的能量为何与 $l$ 无关这一表面上的巧合.
\end{enumerate}
注意,问题的旋转对称性可以解释式~\eqref{eq:hydrogen-bound-state} 中的能量与调和多项式 $q$ 的选择无关.然而,旋转对称性无法解释不同 $l$ 值(从而波函数具有不同径向依赖)对应的状态为何具有相同能量.这个表面上的巧合将由问题的另一个对称性解释,而该对称性可以用 Runge--Lenz 向量表示.另见文献 [17] 第 7 节,其中给出了关于氢原子能级及其重数的另一种(但相关的)解释.

涉及 Runge--Lenz 向量的若干计算虽然初等,却相当繁琐.我们把这些计算推迟到第~\ref{sec:runge-lenz-calculations} 节.

\subsection{一些记号}
\label{subsec:runge-lenz-notation}

为使记号尽可能简单,本节采用 Einstein 求和约定:重复出现的指标总表示求和,即使没有写出求和号.本节中求和范围总是 $1$ 到 $3$.例如,利用这个约定,$\mathbb R^3$ 中两个向量 $u,v$ 的点积写作
\[
 u\cdot v=u_jv_j,
\]
这样就不必显式写出指标 $j$ 的求和.

我们还会频繁使用全反对称符号 $\varepsilon_{jkl}$,其中 $j,k,l$ 取 $1,2,3$,定义如下.

\begin{definition}[Levi--Civita 符号]\label{def:levi-civita-symbol}
对 $j,k,l\in\{1,2,3\}$,定义 $\varepsilon_{jkl}$ 为
\[
 \varepsilon_{jkl}=\begin{cases}
 1,&(j,k,l)\text{ 是 }(1,2,3)\text{ 的偶置换},\\
 -1,&(j,k,l)\text{ 是 }(1,2,3)\text{ 的奇置换},\\
 0,&j,k,l\text{ 中有两个相等}.
 \end{cases}
\]
\end{definition}

例如,$\varepsilon_{321}=-1$,$\varepsilon_{212}=0$.利用求和约定,$\mathfrak{so}(3)$ 基 $\{F_1,F_2,F_3\}$ 的交换关系可以写成
\[
 [F_j,F_k]=\varepsilon_{jkl}F_l.
\label{eq:so3-epsilon-commutator}
\]
例如,取 $j=1,k=2$,关于 $l$ 的求和中只有 $l=3$ 的项非零,于是恢复 $[F_1,F_2]=F_3$.

\subsection{重新考察经典 Runge--Lenz 向量}
\label{subsec:classical-runge-lenz-revisited}

在第~\ref{sec:kepler-problem} 节中,我们已经在 $1/r$ 势下运动的粒子的经典力学中引入了 Runge--Lenz 向量 $A$.在转向量子版本之前,还需要得到 $A$ 的几个性质.考虑 $\mathbb R^3$ 中 Hamilton 函数为
\[
 H(x,p)=\frac{|p|^2}{2\mu}-\frac{Q^2}{|x|}.
\label{eq:classical-hydrogen-hamiltonian}
\]
这就是经典 Kepler 问题的 Hamilton 函数,只是把行星的质量 $m$ 换成电子--质子系统的约化质量 $\mu$,并把常数 $k:=mMG$ 换成 $Q^2$.

对于式~\eqref{eq:classical-hydrogen-hamiltonian} 中的 Hamilton 函数,Runge--Lenz 向量为
\[
 A(x,p)=\frac1{\mu Q^2}p\times J-\frac{x}{|x|},
 \qquad J:=x\times p.
\label{eq:classical-runge-lenz-vector}
\]
根据命题~\ref{pro:runge-lenz-conservation},Runge--Lenz 向量是经典 Kepler 问题中的守恒量;此外,任意径向势中的 $H$ 与 $J$ 也都是守恒量.由第~\ref{sec:kepler-problem} 节的结果,有
\[
 A\cdot J=0,
\qquad
 |A|^2=1+\frac{2H}{\mu Q^4}|J|^2.
\label{eq:classical-runge-lenz-algebraic-relations}
\]

\begin{lemma}[经典 Runge--Lenz 向量的 Poisson 括号]\label{lem:classical-runge-lenz-poisson}
式~\eqref{eq:classical-hydrogen-hamiltonian} 中的 Runge--Lenz 向量 $A$ 与 Hamilton 函数 $H$ 满足
\[
 \{A_j,H\}=0,
 \qquad
 \{A_j,A_m\}=-\frac2{\mu Q^4}\varepsilon_{jml}J_lH.
\label{eq:classical-runge-lenz-poisson}
\]
\end{lemma}

Runge--Lenz 向量是守恒量这一点已经证明过(命题~\ref{pro:runge-lenz-conservation}),而根据命题~\ref{pro:derivative-poisson},这等价于 $\{A_j,H\}=0$.式~\eqref{eq:classical-runge-lenz-poisson} 的第二个关系推迟到第~\ref{sec:runge-lenz-calculations} 节证明.

下面引入 Runge--Lenz 向量、角动量和 Hamilton 函数的某些组合,它们在 Poisson 括号下构成一个 Lie 代数.这些函数的构造需要对 Hamilton 函数取平方根,因此必须把相空间分为正能量部分和负能量部分.我们主要关心负能量情形.

\begin{definition}[负能量下的归一化 Runge--Lenz 向量]\label{def:classical-normalized-runge-lenz}
令 $U^-$ 表示经典相空间的负能量部分:
\[
 U^-:=\{(x,p)\in\mathbb R^6\mid H(x,p)<0\}.
\]
在 $U^-$ 上定义归一化的 Runge--Lenz 向量
\[
 B=\sqrt{\frac{\mu Q^4}{2|H|}}\,A.
\]
并在 $U^-$ 上定义向量值函数 $I,K$:
\[
 I=\frac{J+B}{2},
 \qquad
 K=\frac{J-B}{2}.
\]
\end{definition}

\begin{theorem}[经典 Runge--Lenz 向量产生的 $\mathfrak{so}(4)$ 对称性]\label{thm:classical-so4-symmetry}
函数 $I,K$ 与 Hamilton 函数 Poisson 对易,并且在负能量集合 $U^-$ 上满足
\[
 \{I_j,I_k\}=\varepsilon_{jkl}I_l,
 \qquad
 \{K_j,K_k\}=\varepsilon_{jkl}K_l,
 \qquad
 \{I_j,K_k\}=0.
\label{eq:classical-runge-lenz-so4-brackets}
\]
此外,$I,K$ 满足代数关系
\[
 |I|^2=|K|^2=\frac{\mu Q^4}{8|H|}.
\label{eq:classical-runge-lenz-casimir}
\]
\end{theorem}
}

\def\ChapterEighteenSOFourSpin{
\subsection{\(\mathfrak{so}(4)\) 的表示}
\label{subsec:so4-representations}

根据定理~\ref{thm:quantum-so4-symmetry} 中的交换关系,在负能量子空间 $V^-$ 上可以定义 Lie 代数
\[
 \mathfrak{so}(4)\cong\mathfrak{so}(3)\oplus\mathfrak{so}(3)
\]
的表示 $\pi$,令
\[
 \pi(F_j,0)=\frac1{i\hbar}\widehat I_j,
 \qquad
 \pi(0,F_j)=\frac1{i\hbar}\widehat K_j.
\label{eq:quantum-so4-representation}
\]
因此需要分类 $\mathfrak{so}(3)\oplus\mathfrak{so}(3)$ 的不可约有限维表示;下面的命题给出这个分类.

\begin{proposition}[\(\mathfrak{so}(3)\oplus\mathfrak{so}(3)\) 的不可约表示]\label{prop:so4-irreducible-representations}
设 $V_k,V_l$ 分别是维数为 $2k+1$ 和 $2l+1$ 的不可约 $\mathfrak{so}(3)$ 表示.按照注记~\ref{rem:tensor-product-direct-product} 中的方式把 $V_k\otimes V_l$ 看作 $\mathfrak{so}(3)\oplus\mathfrak{so}(3)$ 的表示,则它不可约.此外,$\mathfrak{so}(3)\oplus\mathfrak{so}(3)$ 的每个不可约有限维表示都同构于某个 $V_k\otimes V_l$,并且有序对 $(k,l)$ 唯一确定.

对任意表示 $V_k\otimes V_l$,定义 Casimir 算子
\[
 C_1=\sum_{j=1}^3\pi_k(F_j)^2\otimes I,
 \qquad
 C_2=\sum_{j=1}^3I\otimes\pi_l(F_j)^2.
\]
则
\[
 C_1=-k(k+1)I,
 \qquad
 C_2=-l(l+1)I.
\]
\end{proposition}

\begin{proof}
要分类 $\mathfrak{so}(3)\oplus\mathfrak{so}(3)$ 的不可约表示,可以直接引用 Lie 代数直和表示的一般理论.不过,也可以用分类 $\mathfrak{so}(3)$ 不可约表示时所采用的同类论证给出直接证明;这里略去计算细节.Casimir 算子的结论由命题~\ref{prop:so3-casimir} 立即得到.
\end{proof}

在 $V^-$ 的任意有限维子空间中,若该子空间在式~\eqref{eq:quantum-so4-representation} 的 $\mathfrak{so}(3)\oplus\mathfrak{so}(3)$ 作用下不变且不可约,则两个 Casimir 算子分别为
\[
 C_1=-\frac{\widehat I\cdot\widehat I}{\hbar^2},
 \qquad
 C_2=-\frac{\widehat K\cdot\widehat K}{\hbar^2}.
\]
由于定理~\ref{thm:quantum-so4-symmetry} 给出在 $V^-$ 上 $\widehat I\cdot\widehat I=\widehat K\cdot\widehat K$,所以 $V^-$ 中出现的所有 $\mathfrak{so}(3)\oplus\mathfrak{so}(3)$ 不可约表示都具有 $V_k\otimes V_k$ 的形式.

\begin{theorem}[氢原子能级上的 \(\mathfrak{so}(4)\) 表示]\label{thm:hydrogen-so4-eigenspace}
令 $W^{(n)}$ 表示 Hamilton 算子特征值为 $E_n$ 的特征子空间.则 $W^{(n)}$ 在式~\eqref{eq:quantum-so4-representation} 的 $\mathfrak{so}(3)\oplus\mathfrak{so}(3)$ 作用下不变且不可约.更具体地,
\[
 W^{(n)}\cong V_k\otimes V_k,
\]
作为 $\mathfrak{so}(3)\oplus\mathfrak{so}(3)$ 的表示,其中 $k=(n-1)/2$,$V_k$ 是维数 $2k+1=n$ 的 $\mathfrak{so}(3)$ 不可约表示.
\end{theorem}

\begin{corollary}[Runge--Lenz 向量给出的氢原子能级]\label{cor:hydrogen-runge-lenz-energy}
若 $n,k,W^{(n)}$ 如定理~\ref{thm:hydrogen-so4-eigenspace},则对所有 $\psi\in W^{(n)}$,
\[
 \widehat I\cdot\widehat I\,\psi
 =\widehat J\cdot\widehat J\,\psi
 =\hbar^2k(k+1)\psi.
\]
结合式~\eqref{eq:quantum-runge-lenz-so4-casimir},可以解出 $\widehat H$ 在 $W^{(n)}$ 上的特征值:
\[
 E_n=-\frac{\mu Q^4}{8\hbar^2\left(k+\frac12\right)^2}
 =-\frac{\mu Q^4}{2\hbar^2n^2}.
\]
\end{corollary}

推论~\ref{cor:hydrogen-runge-lenz-energy} 中的 $E_n$ 表达式与定理~\ref{thm:hydrogen-bound-states} 中相同.推论的证明最令人惊讶之处在于,它纯粹是代数性的:只依赖于算子 $\widehat I_k,\widehat K_l$ 之间的交换关系,以及式~\eqref{eq:quantum-runge-lenz-so4-casimir} 中 Hamilton 算子 $\widehat H$ 与这些算子之间的关系.

\begin{proof}[推论~\ref{cor:hydrogen-runge-lenz-energy} 的证明]
容易看出,当限制在 $\mathfrak{so}(3)\oplus\mathfrak{so}(3)$ 作用下的不可约子空间上时,$\widehat I\cdot\widehat I$ 与 $\widehat K\cdot\widehat K$ 分别等于 $-\hbar^2C_1$ 与 $-\hbar^2C_2$,其中 $C_1,C_2$ 是命题~\ref{prop:so4-irreducible-representations} 中的 Casimir 算子.因此,如果 $W^{(n)}\cong V_k\otimes V_k$,其中 $k=(n-1)/2$,那么 $\widehat I\cdot\widehat I$ 与 $\widehat K\cdot\widehat K$ 都等于 $\hbar^2k(k+1)I$.另一方面,式~\eqref{eq:quantum-runge-lenz-so4-casimir} 把这两个算子与 $\widehat H$ 联系起来,由此即可解出 $E_n$.
\end{proof}

\begin{proof}[定理~\ref{thm:hydrogen-so4-eigenspace} 的证明]
由于 $\widehat A$ 和 $\widehat J$ 的每个分量都与 $\widehat H$ 对易,$\widehat I$ 和 $\widehat K$ 的每个分量也都与 $\widehat H$ 对易.因此 $\widehat H$ 的每个特征子空间在 $\widehat I$ 与 $\widehat K$ 的作用下不变.由于 $\widehat I$ 与 $\widehat K$ 自伴且 $W^{(n)}$ 有限维,$W^{(n)}$ 分解为不可约不变子空间的直和.由命题~\ref{prop:so4-irreducible-representations},这些不可约子空间具有 $V_k\otimes V_l$ 的形式,其中 $V_k,V_l$ 分别是维数为 $2k+1$ 和 $2l+1$ 的不可约表示.

然而,限制在 $W^{(n)}$ 的某个不可约子空间上时,$\widehat I\cdot\widehat I$ 与 $\widehat K\cdot\widehat K$ 分别等于 $-\hbar^2C_1$ 和 $-\hbar^2C_2$.由于在整个 $V^-$ 上 $\widehat I\cdot\widehat I=\widehat K\cdot\widehat K$,所以每个不可约子空间中 $C_1,C_2$ 的特征值必须相等.因此 $k=l$,即只会出现 $V_k\otimes V_k$ 形式的不可约子空间.

在某个 $W^{(n)}$ 的不可约子空间与 $V_k\otimes V_k$ 的同构下,$\widehat I_k$ 与 $\widehat K_k$ 分别作用为 $i\hbar F_k\otimes I$ 与 $i\hbar I\otimes F_k$,其中 $F_k$ 是 $\mathfrak{so}(3)$ 的通常基.由于 $\widehat J=\widehat I+\widehat K$,每个 $\widehat J_k$ 作用为
\[
 i\hbar(F_k\otimes I+I\otimes F_k).
\]
这意味着,在 $\widehat J_k$ 的作用下,$V_k\otimes V_k$ 可以看作两个 $\mathfrak{so}(3)$ 表示的张量积,并按照定义~\ref{def:representation-sum-tensor-dual} 的方式看作另一个 $\mathfrak{so}(3)$ 表示.于是由命题~\ref{prop:so3-tensor-product-decomposition},
\[
 V_k\otimes V_k\cong V_0\oplus V_1\oplus\cdots\oplus V_{2k}.
\label{eq:hydrogen-so4-tensor-decomposition}
\]

另一方面,由定理~\ref{thm:hydrogen-bound-states},$W^{(n)}$ 在 $\mathfrak{so}(3)$ 作用下分解为
\[
 V_0\oplus V_1\oplus\cdots\oplus V_{n-1}.
\label{eq:hydrogen-angular-decomposition}
\]
因此,形如 $V_k\otimes V_k$ 的空间必须就是整个 $W^{(n)}$;否则平凡表示 $V_0$ 会在 $W^{(n)}$ 中出现不止一次.比较式~\eqref{eq:hydrogen-so4-tensor-decomposition} 与式~\eqref{eq:hydrogen-angular-decomposition},得到 $2k=n-1$,即定理中的结论.
\end{proof}

定理~\ref{thm:hydrogen-so4-eigenspace} 的证明在一定程度上依赖第~\ref{sec:hydrogen-bound-states} 节的结果.不过,仅利用 Runge--Lenz 向量的代数运算,我们仍可以论证 $\widehat H$ 的特征值必须具有推论~\ref{cor:hydrogen-runge-lenz-energy} 给出的形式.但这样我们无法知道对每个正整数 $n$,$E_n$ 确实是 $\widehat H$ 的特征值;也无法知道每个特征子空间 $W^{(n)}$ 在 $\mathfrak{so}(4)$ 的作用下不可约.仅凭代数关系,$W^{(n)}$ 甚至可能具有例如 $2n^2$ 而非 $n^2$ 的维数.

\section{自旋的作用}
\label{sec:spin-role-hydrogen}

电子的自旋是 $1/2$.如第~\ref{sec:spin} 节所述,这意味着电子的 Hilbert 空间为
\[
 L^2(\mathbb R^3)\widehat\otimes V_{1/2},
\]
其中 $V_{1/2}$ 是一个二维向量空间,承载 $SO(3)$ 的不可约射影酉表示.到目前为止,我们在计算中忽略了自旋,原因很简单:在第一近似下,自旋不影响计算.具体而言,带自旋氢原子最简单的模型中,Hamilton 算子只是 $\widehat H\otimes I$,其中 $\widehat H$ 是作用于 $L^2(\mathbb R^3)$ 的式~\eqref{eq:hydrogen-hamiltonian} 中的算子.对任意 $n>0$,取
\[
 \psi_{n,l,m}\otimes e_j,
\]
其中 $\psi_{n,l,m}$ 如式~\eqref{eq:hydrogen-spherical-eigenfunction},而 $\{e_1,e_2\}$ 是 $V_{1/2}$ 的一组基,就得到 $\widehat H\otimes I$ 特征值为 $E_n$ 的一组特征向量.

不过,从旋转对称性的角度看,基 $\psi_{n,l,m}\otimes e_j$ 并不是最自然的基.我们应当按照 $SO(3)$ 的射影作用,把特征子空间分解为不可约不变子空间,其中 $SO(3)$ 同时作用于 $L^2(\mathbb R^3)$ 与 $V_{1/2}$.我们已经把 $L^2(\mathbb R^3)$ 中的特征子空间分解为不可约不变子空间,即固定 $n,l$ 而让 $m$ 变化时 $\psi_{n,l,m}$ 的张成空间.因此,要得到 $L^2(\mathbb R^3)\widehat\otimes V_{1/2}$ 中的不可约不变子空间,可以使用第~\ref{sec:tensor-products-angular-momentum} 节的“角动量相加”方法.根据命题~\ref{prop:so3-tensor-half},当 $l=0$ 时 $V_l\otimes V_{1/2}$ 不可约,当 $l>0$ 时它同构于 $V_{l+1/2}\oplus V_{l-1/2}$.

例如,考虑 $n=3,l=1$ 的情形,即传统化学术语中的“$3p$ 态”.由于 $V_1\otimes V_{1/2}$ 分解为 $V_{3/2}\oplus V_{1/2}$,考虑自旋后得到一个四维空间和一个二维空间.沿着命题~\ref{prop:so3-tensor-half} 的证明,可以得到这两个空间的基.

上一段所述的分解对于考察氢原子的“精细结构”至关重要.我们使用 Hamilton 算子~\eqref{eq:hydrogen-hamiltonian} 的氢原子模型只是一阶近似.更现实的模型会纳入各种修正,包括辐射修正、原子核的有限大小和“自旋--轨道耦合”等.自旋--轨道耦合会在 Hamilton 算子中加入含有算子 $\widehat J\cdot\sigma$ 的项,其中 $\sigma_1,\sigma_2,\sigma_3$ 描述 $\mathfrak{so}(3)$ 在 $V_{1/2}$ 上的作用.加入这项后,Hamilton 算子不再具有某个作用于 $L^2(\mathbb R^3)$ 的算子 $A$ 所给出的 $A\otimes I$ 形式.因此,我们不能把自旋简单地附加到计算末尾,而必须从一开始就将它纳入考虑.

氢原子 Hamilton 算子的各种修正会降低能级的重数.比如,几乎任何修正都会破坏给定 $n$ 时能级与 $l$ 无关这一性质,因为 Hamilton 算子中的修正项一般不与量子 Runge--Lenz 向量对易.然而,构成氢原子精细结构的所有修正都保持问题的旋转对称性.因此,简单模型中出现的同样的 $SO(3)$ 不可约表示,在加入修正后仍会出现.例如当 $n=2,l=1$ 时,我们仍会有一个四维空间和一个二维空间,但这两个空间不再具有相同的能量.

}

\def\ChapterEighteenComputations{
\section{Runge--Lenz 计算}
\label{sec:runge-lenz-calculations}

本节补充第~\ref{sec:quantum-runge-lenz} 节中略去的许多计算.原则上,所有计算都是初等的,但其中有一些并不明显的技巧可以简化代数运算.我们会频繁使用“按向量变换的函数”(经典一侧)和“向量算子”(量子一侧)的概念,包括第~\ref{sec:vectors-vector-operators} 节的命题~\ref{prop:classical-vector-poisson-relations} 与命题~\ref{prop:vector-operator-commutators}.特别地,位置 $x$、动量 $p$、角动量 $j$ 和 Runge--Lenz 向量 $A$ 都按向量变换,相应的量子量都是向量算子.(比较习题 7.)用第~\ref{subsec:runge-lenz-notation} 节的“$\varepsilon$”记号,命题~\ref{prop:vector-operator-commutators} 变为
\[
 \frac1{i\hbar}[C_j,\widehat J_k]
 =\frac1{i\hbar}[\widehat J_j,C_k]
 =\varepsilon_{jkl}C_l.
\label{eq:vector-operator-epsilon-commutator}
\]

量子力学计算中会出现若干“量子修正”:向量算子的点积与叉积不再完全遵循经典情形中的规则.

\begin{lemma}[Levi--Civita 符号的恒等式]\label{lem:epsilon-identities}
定义~\ref{def:levi-civita-symbol} 中的 $\varepsilon$ 符号满足
\[
 \varepsilon_{jkl}\varepsilon_{jmn}=\delta_{km}\delta_{ln}-\delta_{kn}\delta_{lm},
 \qquad
 \varepsilon_{jkl}\varepsilon_{jkm}=2\delta_{lm}.
\]
\end{lemma}

这些结果并不难证明,留给读者(习题 6).下面的向量算子叉积恒等式将会有用.

\begin{lemma}[向量算子叉积恒等式]\label{lem:vector-operator-cross-product}
若 $C,D,E$ 是任意向量算子,则
\[
 C\cdot(D\times E)=(C\times D)\cdot E,
\label{eq:operator-cross-dot-cyclic}
\]
\[
 C\times D+D\times C=\varepsilon_{jkl}[C_k,D_l],
\label{eq:operator-cross-symmetrized}
\]
\[
 C\times C=\frac12\varepsilon_{jkl}[C_k,C_l].
\label{eq:operator-self-cross}
\]
特别地,如果 $C$ 的不同分量彼此对易,则 $C\times C=0$.最后,
\[
 (C\times(D\times E))_j=C_kD_jE_k-C_kD_kE_j.
\label{eq:operator-triple-cross}
\]
\end{lemma}

这些结果的特殊情形为
\[
 \widehat J\times P+P\times\widehat J=2i\hbar P,
\label{eq:angular-momentum-momentum-cross}
\]
\[
 \widehat J\times\widehat J=i\hbar\widehat J.
\label{eq:angular-momentum-self-cross}
\]
注意,如果 $D,E$ 的各分量彼此对易,那么式~\eqref{eq:operator-triple-cross} 的右端化为经典表达式 $(C\cdot E)D-(C\cdot D)E$.利用式~\eqref{eq:angular-momentum-momentum-cross},可以很容易验证式~\eqref{eq:quantum-runge-lenz-alternative} 中量子 Runge--Lenz 向量的另一种表达式.

\begin{proof}
式~\eqref{eq:operator-cross-dot-cyclic} 的右端按分量计算为 $\varepsilon_{jkl}C_kD_lE_j$.注意 $\varepsilon_{jkl}=\varepsilon_{klj}$,再重新标记指标,就得到 $\varepsilon_{jkl}C_jD_kE_l$,这等于左端.

对于式~\eqref{eq:operator-cross-symmetrized},计算得
\[
\begin{aligned}
 (C\times D+D\times C)_j
 &=\varepsilon_{jkl}C_kD_l+\varepsilon_{jkl}D_kC_l\\
 &=\varepsilon_{jkl}C_kD_l+\varepsilon_{jkl}C_lD_k
 -\varepsilon_{jkl}[C_l,D_k].
\end{aligned}
\label{eq:operator-cross-symmetrized-calculation}
\]
由于 $\varepsilon_{jkl}=-\varepsilon_{jlk}$,交换指标 $k,l$ 可得 $\varepsilon_{jkl}C_lD_k=-\varepsilon_{jkl}C_kD_l$,所以第二行中的前两项相消.剩下的项通过重新标记指标即可写成所声称的形式.式~\eqref{eq:operator-self-cross} 是式~\eqref{eq:operator-cross-symmetrized} 取 $D=C$ 的特殊情形.最后,式~\eqref{eq:operator-triple-cross} 直接由引理~\ref{lem:epsilon-identities} 得到.

式~\eqref{eq:angular-momentum-momentum-cross} 与式~\eqref{eq:angular-momentum-self-cross} 分别由式~\eqref{eq:operator-cross-symmetrized} 和式~\eqref{eq:operator-self-cross} 得到,因为 $\widehat J$ 与 $P$ 都是向量算子.引理~\ref{lem:epsilon-identities} 给出所需的最后一步.
\end{proof}

现在证明第~\ref{sec:quantum-runge-lenz} 节的结果.我们只证明量子版本,因为经典版本极为相似,只需忽略某些量子修正.

\begin{proof}[引理~\ref{lem:quantum-runge-lenz-commutator} 的第一部分]
先证明每个 $\widehat A_j$ 与 $\widehat H$ 对易.由于 $\widehat H$ 与 $\widehat J$ 对易,有
\[
 [\widehat A_j,\widehat H]
 =\frac1{\mu Q^2}\varepsilon_{jkl}\left(
 \frac12[ P_k,\widehat H]\widehat J_l
 -\frac12\widehat J_k[P_l,\widehat H]
 -\left[\frac{X_j}{|X|},\widehat H\right]\right).
\label{eq:quantum-runge-lenz-hamiltonian-commutator-start}
\]
另一方面,由于各个 $P$ 分量彼此对易,
\[
 [P_k,\widehat H]=-Q^2\left[P_k,\frac1{|X|}\right]
 =-i\hbar Q^2\frac{X_k}{|X|^3}.
\]
因此
\[
\begin{aligned}
 \varepsilon_{jkl}[P_k,\widehat H]\widehat J_l
 &=-i\hbar Q^2\varepsilon_{jkl}\varepsilon_{lmn}
 \frac{X_k}{|X|^3}X_mP_n\\
 &=-i\hbar Q^2(\delta_{jm}\delta_{kn}-\delta_{jn}\delta_{km})
 \frac{X_k}{|X|^3}X_mP_n\\
 &=-i\hbar Q^2\frac1{|X|^3}
 (X_nX_jP_n-X_mX_mP_j)\\
 &=-i\hbar Q^2\frac1{|X|^3}
 \bigl(X_j(X\cdot P)-(X\cdot X)P_j\bigr).
\end{aligned}
\label{eq:quantum-runge-lenz-commutator-term1}
\]
类似地计算 $\varepsilon_{jkl}\widehat J_k[P_l,\widehat H]$.注意
\[
 \widehat J_k=\varepsilon_{kmn}X_mP_n=\varepsilon_{kmn}P_nX_m,
\]
因为当 $m=n$ 时 $\varepsilon_{kmn}=0$,而除这种情形外 $X_m,P_n$ 对易.于是
\[
 \varepsilon_{jkl}\widehat J_k[P_l,\widehat H]
 =-i\hbar Q^2\frac1{|X|^3}
 \bigl(P_j(X\cdot X)-(P\cdot X)X_j\bigr).
\label{eq:quantum-runge-lenz-commutator-term2}
\]

另一方面,由于各 $X$ 分量彼此对易,
\[
\begin{aligned}
 \left[\frac{X_j}{|X|},\widehat H\right]
 &=\left[\frac{X_j}{|X|},\frac{P^2}{2\mu}\right]\\
 &=\frac1{2\mu}\left[\frac{X_j}{|X|},P_kP_k\right]\\
 &=\frac1{2\mu}\left(\left[\frac{X_j}{|X|},P_k\right]P_k
 +P_k\left[\frac{X_j}{|X|},P_k\right]\right)\\
 &=\frac{i\hbar}{2\mu}\left(\delta_{jk}\frac1{|X|}-\frac{X_jX_k}{|X|^3}\right)P_k
 +\frac{i\hbar}{2\mu}P_k\left(\delta_{jk}\frac1{|X|}-\frac{X_jX_k}{|X|^3}\right)\\
 &=\frac{i\hbar}{2\mu}\left(P_j-\frac{X_j}{|X|^3}(X\cdot P)\right)
 +\frac{i\hbar}{2\mu}\left(P_j-(P\cdot X)\frac{X_j}{|X|^3}\right).
\end{aligned}
\label{eq:quantum-runge-lenz-position-commutator}
\]
将式~\eqref{eq:quantum-runge-lenz-commutator-term1}、式~\eqref{eq:quantum-runge-lenz-commutator-term2} 和式~\eqref{eq:quantum-runge-lenz-position-commutator} 代入式~\eqref{eq:quantum-runge-lenz-hamiltonian-commutator-start},各项正好相消.例如,式~\eqref{eq:quantum-runge-lenz-commutator-term1} 中含有 $(X_j/|X|^3)(X\cdot P)$ 的项,而式~\eqref{eq:quantum-runge-lenz-position-commutator} 中有相应的抵消项.这就证明了 $[\widehat A_j,\widehat H]=0$.
\end{proof}

在继续研究 Runge--Lenz 向量的其他性质之前,先验证一些稍后要用到的结果.与相应的经典结果相比,这里会出现量子修正.

\begin{lemma}[Runge--Lenz 计算所需的正交关系]\label{lem:runge-lenz-orthogonality}
与经典情形一样,向量算子满足“正交”关系
\[
 \widehat J\cdot P=P\cdot\widehat J=0,
\label{eq:runge-lenz-jp-orthogonal}
\]
\[
 \widehat J\cdot X=X\cdot\widehat J=0,
\label{eq:runge-lenz-jx-orthogonal}
\]
\[
 (P\times\widehat J)\cdot\widehat J
 =\widehat J\cdot(P\times\widehat J)=0.
\label{eq:runge-lenz-pj-j-orthogonal}
\]
不过,$P$ 与 $P\times\widehat J$ 的点积存在量子修正:
\[
 P\cdot(P\times\widehat J)=0,
\label{eq:runge-lenz-p-pj}
\]
\[
 (P\times\widehat J)\cdot P=2i\hbar(P\cdot P).
\label{eq:runge-lenz-pj-p}
\]
最后还有
\[
 (P\times\widehat J)\cdot(P\times\widehat J)=(P\cdot P)(\widehat J\cdot\widehat J),
\label{eq:runge-lenz-pj-square}
\]
\[
 X\cdot(P\times\widehat J)=\widehat J\cdot\widehat J,
\label{eq:runge-lenz-x-pj}
\]
\[
 (P\times\widehat J)\cdot X=\widehat J\cdot\widehat J+2i\hbar P\cdot X.
\label{eq:runge-lenz-pj-x}
\]
\end{lemma}

\begin{proof}
由式~\eqref{eq:operator-cross-dot-cyclic} 和式~\eqref{eq:operator-self-cross},
\[
 \widehat J\cdot P=(X\times P)\cdot P=X\cdot(P\times P)=0,
\]
因为 $P$ 的不同分量彼此对易.相同的论证说明 $P\cdot\widehat J$、$\widehat J\cdot X$ 和 $X\cdot\widehat J$ 都为零.

计算 $(P\times\widehat J)\cdot\widehat J$ 时,先使用式~\eqref{eq:operator-cross-dot-cyclic},再使用式~\eqref{eq:angular-momentum-self-cross},最后使用 $P\cdot\widehat J=0$.对于 $\widehat J\cdot(P\times\widehat J)$,利用式~\eqref{eq:angular-momentum-momentum-cross} 把 $P\times\widehat J$ 改写成 $\widehat J\times P$.修正项含有 $P$,它与 $\widehat J$ 的点积为零,因此结果同样为零.

再次使用式~\eqref{eq:operator-cross-dot-cyclic} 和式~\eqref{eq:operator-self-cross} 可得式~\eqref{eq:runge-lenz-p-pj}.要得到式~\eqref{eq:runge-lenz-pj-p},先利用式~\eqref{eq:angular-momentum-momentum-cross} 把 $P\times\widehat J$ 改写成 $\widehat J\times P$,再应用式~\eqref{eq:runge-lenz-p-pj}.

为证明式~\eqref{eq:runge-lenz-pj-square},应用式~\eqref{eq:operator-cross-dot-cyclic} 与式~\eqref{eq:operator-triple-cross},得到
\[
 (P\times\widehat J)\cdot(P\times\widehat J)
 =P_j\widehat J_kP_j\widehat J_k-P_j\widehat J_kP_k\widehat J_j.
\label{eq:runge-lenz-pj-square-expanded}
\]
右端第二项为零,因为 $\widehat J\cdot P=0$.对于第一项,把 $\widehat J_k$ 向右移过 $P_j$;这会产生所需项以及修正项 $i\hbar\varepsilon_{kjl}P_jP_l\widehat J_k$.该修正项为零,因为 $P_j,P_l$ 对易,而 $\varepsilon_{kjl}$ 在交换 $j,l$ 时变号.于是得到式~\eqref{eq:runge-lenz-pj-square}.式~\eqref{eq:runge-lenz-x-pj} 直接由式~\eqref{eq:operator-cross-dot-cyclic} 和角动量定义得到;式~\eqref{eq:runge-lenz-pj-x} 由式~\eqref{eq:operator-cross-dot-cyclic} 与式~\eqref{eq:operator-cross-symmetrized} 得到.
\end{proof}

定理~\ref{thm:classical-so4-symmetry} 使用了上一小节引入的求和约定.该定理的证明初等但相当繁琐,推迟到第~\ref{sec:runge-lenz-calculations} 节.

在 $U^-$ 上,函数 $I_1,I_2,I_3,K_1,K_2,K_3$ 的张成空间与 $B_1,B_2,B_3,J_1,J_2,J_3$ 的张成空间相同,并且在 Poisson 括号下构成一个六维 Lie 代数.比较 $I$ 之间及 $K$ 之间的 Poisson 括号关系与 $\mathfrak{so}(3)$ 基 $F_1,F_2,F_3$ 的关系,可知 $I$ 的张成空间和 $K$ 的张成空间都同构于 $\mathfrak{so}(3)$(或者说同构于 $\mathfrak{su}(2)$).又因为每个 $I_j$ 都与每个 $K_k$ 对易,所以由 $I$ 与 $K$ 张成的六维 Lie 代数同构于 $\mathfrak{so}(3)\oplus\mathfrak{so}(3)$.同时,正如习题 4 所示,$\mathfrak{so}(3)\oplus\mathfrak{so}(3)$ 同构于 Lie 代数 $\mathfrak{so}(4)$.由于所有 $I_j,K_j$ 都与 Hamilton 函数 Poisson 对易,我们说 Kepler 问题具有 $\mathfrak{so}(4)$ 对称性.这与在一般径向势产生的力作用下运动的 $\mathbb R^3$ 粒子不同;后者只有 $\mathfrak{so}(3)$ 对称性.

更精确地说,“$\mathfrak{so}(4)$ 对称性”只在经典相空间的负能量子集 $U^-$ 上成立.在正能量子集 $U^+$ 上,函数 $B_1,B_2,B_3,J_1,J_2,J_3$ 仍然张成一个六维 Lie 代数,但这个 Lie 代数不再同构于 $\mathfrak{so}(4)$,而是同构于 $\mathfrak{so}(3,1)$;这里 $\mathfrak{so}(3,1)$ 是保持二次型 $x_1^2+x_2^2+x_3^2-x_4^2$ 的 $4\times4$ 矩阵群的 Lie 代数.$U^+$ 上的公式不同,是因为相关 Poisson 括号计算包含函数 $H/|H|$:它在 $U^+$ 上取值 $1$,在 $U^-$ 上取值 $-1$.(其中一个 $\sqrt{|H|}$ 因子来自引理~\ref{lem:classical-runge-lenz-poisson},另一个 $\sqrt{|H|}$ 因子来自 $B$ 的定义.)

}

\def\ChapterEighteenQuantumRungeLenz{
\subsection{量子 Runge--Lenz 向量}
\label{subsec:quantum-runge-lenz-vector}

现在引入经典 Runge--Lenz 向量 $A$ 的量子对应物 $\widehat A$.量子 Runge--Lenz 向量与经典版本具有大部分相同的性质,但会出现几个虽小却关键的“量子修正”.

\begin{definition}[量子 Runge--Lenz 向量]\label{def:quantum-runge-lenz-vector}
定义量子 Runge--Lenz 向量为
\[
 \widehat A=\frac1{\mu Q^2}\left[\frac12( P\times\widehat J-\widehat J\times P)-\frac{X}{|X|}\right].
\label{eq:quantum-runge-lenz-vector}
\]
\end{definition}

注意,在量子情形中,由于因子不对易,$-\widehat J\times P$ 与 $P\times\widehat J$ 并不相同.定义~\ref{def:quantum-runge-lenz-vector} 中采用 $P\times\widehat J$ 与 $\widehat J\times P$ 的特定组合,是因为它给出自伴算子.Runge--Lenz 向量还可以写成
\[
 \widehat A=\frac1{\mu Q^2}\left[P\times\widehat J-i\hbar P-\frac{X}{|X|}\right],
\label{eq:quantum-runge-lenz-alternative}
\]
这一点将在第~\ref{sec:runge-lenz-calculations} 节验证.

为使叙述保持在可控范围内,下面不讨论各种恒等式究竟在哪些精确定义域上成立.

\begin{proposition}[量子 Runge--Lenz 向量的代数关系]\label{prop:quantum-runge-lenz-algebraic}
量子 Runge--Lenz 向量满足
\[
 \widehat A\cdot\widehat J=\widehat J\cdot\widehat A=0,
\]
以及
\[
 \widehat A\cdot\widehat A
 =1+\frac{2\widehat H}{\mu Q^4}\left(\widehat J\cdot\widehat J+\hbar^2\right).
\label{eq:quantum-runge-lenz-square}
\]
\end{proposition}

注意式~\eqref{eq:quantum-runge-lenz-square} 中的“量子修正”:经典表达式 $A\cdot A$ 中的 $J\cdot J$ 被 $\widehat J\cdot\widehat J+\hbar^2$ 替代.这个修正导致式~\eqref{eq:quantum-runge-lenz-so4-casimir} 中的量子修正,而后者对于在推论~\ref{cor:hydrogen-runge-lenz-energy} 中得到正确的能级至关重要.本结果及本节其他结果的证明推迟到第~\ref{sec:runge-lenz-calculations} 节.

\begin{lemma}[量子 Runge--Lenz 向量与 Hamilton 算子的交换关系]\label{lem:quantum-runge-lenz-commutator}
量子 Runge--Lenz 向量 $\widehat A$ 与 Hamilton 算子 $\widehat H$ 满足
\[
 \frac1{i\hbar}[\widehat A_j,\widehat H]=0,
\qquad
 \frac1{i\hbar}[\widehat A_j,\widehat A_m]
 =-\frac2{\mu Q^4}\varepsilon_{jml}\widehat J_l\widehat H.
\label{eq:quantum-runge-lenz-commutator}
\]
\end{lemma}

由于 $\widehat H$ 与旋转对易,它与角动量算子 $\widehat J_l$ 对易.因此在式~\eqref{eq:quantum-runge-lenz-commutator} 中,也可以把 $\widehat J_l\widehat H$ 写成 $\widehat H\widehat J_l$.与经典情形一样,如果把 Runge--Lenz 向量的分量除以 Hamilton 算子的平方根进行归一化,那么这些算子连同角动量算子将构成一个六维 Lie 代数.

\begin{definition}[量子负能量子空间上的归一化算子]\label{def:quantum-normalized-runge-lenz}
令 $V^-$ 表示 $L^2(\mathbb R^3)$ 的负能量子空间,即谱投影 $\mu_{\widehat H}((-\infty,0))$ 的值域.令 $|\widehat H|$ 表示算子 $-\widehat H$ 在 $V^-$ 上的限制.在 $V^-$ 上定义
\[
 \widehat B=\sqrt{\frac{\mu Q^4}{2|\widehat H|}}\,\widehat A.
\]
并像经典情形一样定义
\[
 \widehat I=\frac{\widehat J+\widehat B}{2},
 \qquad
 \widehat K=\frac{\widehat J-\widehat B}{2}.
\]
\end{definition}

任意自伴算子的绝对值都可以用函数演算定义.不过,由于 $\widehat H$ 在 $V^-$ 上的限制按定义是负定的,$|\widehat H|$ 在 $V^-$ 上的限制与 $-\widehat H$ 在 $V^-$ 上的限制相同.算子 $1/|\widehat H|$ 在特征值为 $E_n$ 的能量特征子空间上的限制是 $(1/|E_n|)I$.$\widehat B$ 的各分量是无界算子,定义在 $V^-$ 的适当稠密子空间上.

\begin{theorem}[量子 $\mathfrak{so}(4)$ 对称性]\label{thm:quantum-so4-symmetry}
算子 $\widehat I$ 和 $\widehat K$ 与 Hamilton 算子 $\widehat H$ 对易,并满足
\[
 \frac1{i\hbar}[\widehat I_j,\widehat I_k]=\varepsilon_{jkl}\widehat I_l,
\qquad
 \frac1{i\hbar}[\widehat K_j,\widehat K_k]=\varepsilon_{jkl}\widehat K_l,
\qquad
 \frac1{i\hbar}[\widehat I_j,\widehat K_k]=0.
\label{eq:quantum-runge-lenz-so4-brackets}
\]
此外,它们满足
\[
 \widehat I\cdot\widehat I=\widehat K\cdot\widehat K
 =\frac{\mu Q^4}{8|\widehat H|}-\frac{\hbar^2}{4}.
\label{eq:quantum-runge-lenz-so4-casimir}
\]
\end{theorem}

}

\def\ChapterEighteenRadialRemainder{
尽管定义~\ref{def:radial-representation-subspace} 中对函数的分解在许多方面最方便,但它并不是把球谐函数变成 $\mathbb R^3$ 上函数时通常采用的形式.通常,人们使用极坐标,并考虑
\[
 \psi(r,\theta,\phi)=p(\theta,\phi)g(r),
\]
其中 $p$ 是某个 $V_l$ 元素在球面上的限制.用直角坐标表示,这种分解为
\[
 \psi(x)=p\left(\frac{x}{|x|}\right)g(|x|)=\frac{p(x)}{|x|^l}g(|x|).
\]

\begin{proposition}[球坐标形式下的 Laplace 算子]\label{prop:radial-laplacian-angular}
设 $p\in V_l$,$f$ 是 $(0,\infty)$ 上的光滑函数,并令 $\psi$ 为 $\mathbb R^3\setminus\{0\}$ 上的函数
\[
 \psi(x)=p\left(\frac{x}{|x|}\right)g(|x|).
\]
则对所有 $x\in S^2$ 和 $r\in(0,\infty)$,有
\[
 (\Delta\psi)(rx)=p(x)\left(\frac{d^2g}{dr^2}+\frac2r\frac{dg}{dr}-\frac{l(l+1)}{r^2}g(r)\right).
\label{eq:radial-laplacian-spherical}
\]
\end{proposition}

\begin{proof}
由于 $p$ 是 $l$ 次齐次的,
\[
 p\left(\frac{x}{|x|}\right)=\frac{p(x)}{|x|^l}.
\]
因此
\[
 \psi(x)=p(x)\frac{f(|x|)}{|x|^l}.
\]
应用命题~\ref{prop:radial-laplacian},得到
\[
 \Delta\psi(x)=p(x)\left(\frac{d^2}{dr^2}+\frac{2(l+1)}r\frac{d}{dr}\right)\left(\frac{f(r)}{r^l}\right).
\]
从这里通过直接计算即可验证命题中的公式;计算虽然直截了当,但并无特别启发性.
\end{proof}

在 $\mathbb R^3$ 上写函数还有另一种形式:
\[
 \psi(x)=\frac1{|x|}p\left(\frac{x}{|x|}\right)h(|x|),
\label{eq:radial-reduced-wavefunction}
\]
于是 $h(r)=rg(r)$.在式~\eqref{eq:radial-laplacian-spherical} 中用 $h(r)/r$ 替换 $g(r)$,经过简短计算得到
\[
 (\Delta\psi)(rx)=\frac1{|x|}p(x)\left(\frac{d^2h}{dr^2}-\frac{l(l+1)}{r^2}h(r)\right),
 \qquad x\in S^2.
\label{eq:radial-reduced-laplacian}
\]
采用式~\eqref{eq:radial-reduced-wavefunction} 的形式写波函数很方便,因为对任意径向势都有
\[
 -\frac{\hbar^2}{2m}\Delta\psi+V(|x|)\psi
 =\frac1{|x|}p(x)\left[-\frac{\hbar^2}{2m}\frac{d^2h}{dr^2}+V_{\mathrm{eff}}(r)h(r)\right],
\label{eq:effective-radial-hamiltonian}
\]
其中有效势为
\[
 V_{\mathrm{eff}}(r)=V(r)+\frac{\hbar^2l(l+1)}{2mr^2}.
\label{eq:effective-radial-potential}
\]

注意,式~\eqref{eq:effective-radial-hamiltonian} 中方括号内的量只是普通的一维 Schrödinger 算子,因为式~\eqref{eq:radial-laplacian-spherical} 中的一阶导数项已被消去.尽管式~\eqref{eq:radial-reduced-wavefunction} 很自然,但最终寻找氢原子 Hamilton 算子的束缚态时,式~\eqref{eq:radial-laplacian-function} 的形式最为方便.

正如第~\ref{sec:basic-unbounded-operators} 节后的讨论所说明的,即使 $\psi$ 在 $\mathbb R^3\setminus\{0\}$ 上平方可积,且 $\Delta\psi$ 在 $\mathbb R^3\setminus\{0\}$ 上平方可积,$\psi$ 也未必属于 Laplace 算子的定义域,因为 $\psi$ 的分布 Laplace 算子可能含有一个支集位于原点的项.不过,在氢原子的情形中,我们将考虑式~\eqref{eq:radial-laplacian-function} 形式的函数,其中 $f$ 和 $df/dr$ 在原点附近有界,并且在无穷远处指数衰减.于是命题~\ref{prop:radial-laplacian-domain} 告诉我们,$\psi$ 属于 $\Delta$ 的定义域.

}

\def\ChapterEighteenHydrogen{
\section{氢原子:预备知识}
\label{sec:hydrogen-preliminaries}

氢原子由一个通过带相反电荷粒子之间的电磁吸引力而“束缚”在质子附近的电子组成.氢原子的研究是量子力学中非常重要的检验问题,Schrödinger 方程能够解释氢原子观测到的能级,是量子理论早期的关键成功之一.

质子的质量约为电子质量的 $1800$ 倍.因此,在第一近似下,可以认为质子的位置固定,而电子围绕这个位置“运行”.更精确的分析会考虑质子和电子都绕它们的质心运动.两粒子相对位置的 Hamilton 函数与粒子绕固定中心运动的 Hamilton 函数完全相同,只是把电子质量替换为电子--质子系统的约化质量 $\mu$.(见习题 1.)这里与经典情形中的命题~\ref{pro:reduced-mass} 一样,
\[
 \mu=\frac{m_em_p}{m_e+m_p},
\label{eq:hydrogen-reduced-mass}
\]
其中 $m_e,m_p$ 分别是电子和质子的质量.由于 $m_p\gg m_e$,约化质量与电子质量几乎相同.

分离出质心运动后,电子相对位置的 Hamilton 算子为
\[
 \widehat H=-\frac{\hbar^2}{2\mu}\Delta-\frac{Q^2}{|x|},
\label{eq:hydrogen-hamiltonian}
\]
其中 $Q$ 是电子的电荷.我们使用静电单位制或 Gaussian 单位制等 Coulomb 常数为 $1$ 的单位制.由定理~\ref{thm:kato-rellich},$\widehat H$ 在 $\operatorname{Dom}(\Delta)$ 上自伴,并且有下界.

注意,氢原子的经典 Hamilton 函数 $H(x,p)$ 没有下界.事实上,只需取 $p=0$,并令 $x$ 非常接近原点即可.这种无界性会使假想的经典氢原子表现出奇怪的行为.毕竟,用 $1/r$ 势来模拟氢原子只是一种近似:我们把在粒子位置固定时正确的静电力公式,用在了一个动力学情形中.更现实的氢原子模型会考虑辐射,即带电电子与电磁场之间的相互作用.经典地说,一个绕正电荷原子核运行的负电荷粒子会辐射,从而把能量交给电磁场.经典粒子会迅速向原点盘旋,粒子的能量趋于 $-\infty$,而电磁场的能量趋于 $+\infty$.因此,如果氢原子由经典带电粒子组成,电子会进入“死亡螺旋”,并释放巨大的电磁辐射爆发.

幸运的是,真实粒子并不是这样运动的！实际上,电子是量子粒子.绕质子“运行”的量子电子仍然可以把能量交给电磁场;但量子氢原子的 Hamilton 算子由定理~\ref{thm:kato-rellich} 保证有下界.因此,电子只能下降到基态(最低能量状态),到达基态后便稳定下来.

\section{氢原子的束缚态}
\label{sec:hydrogen-bound-states}

本节的目标是寻找 Hamilton 算子~\eqref{eq:hydrogen-hamiltonian} 的负特征值对应的特征向量.这些特征向量构成“束缚态”,即电子被质子束缚的状态.对每个负数 $E$,考察 $\widehat H$ 特征值为 $E$ 的特征子空间 $V_E$,也就是所有满足 $\widehat H\psi=E\psi$ 的 $\psi\in\operatorname{Dom}(\widehat H)$ 所组成的空间.由于 $\widehat H$ 自伴,因而闭,这个特征子空间是 $L^2(\mathbb R^3)$ 的闭子空间.此外,由于 $\widehat H$ 与旋转对易,$V_E$ 在定义~\ref{def:rotation-representation} 的 $SO(3)$ 在 $L^2(\mathbb R^3)$ 上的自然作用下不变.因此,根据第~\ref{sec:representations-in-l2-r3} 节末尾的讨论,$V_E$ 分解为有限维不可约 $SO(3)$ 不变子空间的直和.

现在寻找 $V_E$ 中的这类子空间.在下面的定理中,假设波函数的径向部分(即定义~\ref{def:radial-representation-subspace} 中的函数 $f$)具有某种非常特殊的形式.分析完这个情形后,我们将说明这样得到的正是 $\widehat H$ 的全部负能量特征向量.

\begin{theorem}[氢原子的束缚态]\label{thm:hydrogen-bound-states}
对每个正整数 $n$,令
\[
 E_n=-\frac{\mu Q^4}{2\hbar^2n^2},
\label{eq:hydrogen-energy-levels}
\]
其中 $Q$ 是电子电荷,$\mu$ 是电子--质子系统的约化质量,并令
\[
 \rho_n(x)=\frac{\sqrt{8\mu|E_n|}}{\hbar}|x|.
\label{eq:hydrogen-rho}
\]
那么,对每个 $l=0,1,\ldots,n-1$,存在一个多项式 $L_{n,l}$,使得对每个次数为 $l$ 的齐次调和多项式 $q$,函数
\[
 \psi(x)=q(x)e^{-\rho_n(x)/2}L_{n,l}(\rho_n(x))
\label{eq:hydrogen-bound-state}
\]
满足
\[
 \widehat H\psi=E_n\psi.
\]
\end{theorem}

由命题~\ref{prop:radial-laplacian-domain},式~\eqref{eq:hydrogen-bound-state} 中的函数属于 $\operatorname{Dom}(\Delta)$,因而由定理~\ref{thm:kato-rellich} 属于 $\operatorname{Dom}(\widehat H)$.多项式 $L_{n,l}$ 就是 Laguerre 多项式.式~\eqref{eq:hydrogen-energy-levels} 中 $-1/n^2$ 的系数是 Rydberg 常数(比较第~\ref{subsec:hydrogen-spectrum} 节).

现在说明定理~\ref{thm:hydrogen-bound-states} 与物理文献中氢原子特征向量的通常表达式之间的联系.首先,物理学家会为调和多项式空间选取一组基 $q_{l,m}$;除归一化常数外,这就是定理~\ref{thm:so3-irreducible-standard-form} 中的基.其次,物理学家使用球坐标写解.换成球坐标时要记住,$q_{l,m}$ 是 $l$ 次齐次的,而 $\rho_n(x)$ 只是到原点距离的某个常数倍.因此得到
\[
 \psi_{n,l,m}(r,\theta,\phi)=Y_{l,m}(\theta,\phi)\rho_n^l e^{-\rho_n/2}L_{n,l}(\rho_n),
\label{eq:hydrogen-spherical-eigenfunction}
\]
其中 $Y_{l,m}(\theta,\phi)$ 是 $q_{l,m}$ 在单位球面上的限制.

\begin{proof}
若 $E$ 是负实数,我们寻找形如 $q(x)f(|x|)$ 的 $\widehat H\psi=E\psi$ 解,其中 $q\in V_l$.如果 $f(r)$ 和 $f'(r)$ 在原点附近有界,命题~\ref{prop:radial-laplacian-domain} 允许我们在 $\mathbb R^3\setminus\{0\}$ 上计算 $\Delta\psi$,而不必担心 $\psi$ 在原点是否可微.由命题~\ref{prop:radial-laplacian},$f$ 满足
\[
 -\frac{\hbar^2}{2\mu}\left(\frac{d^2f}{dr^2}+\frac{2(l+1)}r\frac{df}{dr}\right)-\frac{Q^2}{r}f(r)=Ef(r).
\label{eq:hydrogen-radial-equation}
\]
对大 $r$,含有 $1/r$ 因子的两项可以忽略,于是
\[
 -\frac{\hbar^2}{2\mu}\frac{d^2f}{dr^2}\approx Ef.
\label{eq:hydrogen-radial-asymptotic}
\]
记住 $E<0$,式~\eqref{eq:hydrogen-radial-asymptotic} 告诉我们,在无穷远附近,$f$ 应表现为增长指数与衰减指数的线性组合.由于我们希望得到平方可积解,要求只保留指数衰减项.

因此设解为
\[
 f(r)=\exp\left(-\frac{\sqrt{2\mu|E|}}{\hbar}r\right)g(r),
\label{eq:hydrogen-radial-ansatz}
\]
其中 $g$ 是待定函数.将式~\eqref{eq:hydrogen-radial-ansatz} 代入 $f$ 的方程,等式两边各有相互抵消的 $Eg(r)$ 项,留下
\[
 -\frac{\hbar^2}{2\mu}\left[\frac{d^2g}{dr^2}-2\frac{\sqrt{2\mu|E|}}{\hbar}\frac{dg}{dr}+\frac{2(l+1)}r\frac{dg}{dr}-\frac{2(l+1)}r\frac{\sqrt{2\mu|E|}}{\hbar}g(r)\right]
 =\frac{Q^2}{r}g(r).
\]

引入新变量
\[
 \rho=\frac{\sqrt{8\mu|E|}}{\hbar}r.
\]
作这个变量替换后,方括号中的每一项都得到因子 $\sqrt{8\mu|E|}/\hbar^2$ 的适当幂次;整理后得到
\[
 \rho\frac{d^2g}{d\rho^2}+[2(l+1)-\rho]\frac{dg}{d\rho}
 +\left[\lambda-(l+1)\right]g(\rho)=0,
\label{eq:hydrogen-laguerre-equation}
\]
其中
\[
 \lambda=\frac{Q^2}{\hbar}\sqrt{\frac{\mu}{2|E|}}.
\]
若设 $g(\rho)=\sum_{k=0}^{\infty}a_k\rho^k$,就得到系数的递推关系
\[
 a_{k+1}=a_k\frac{k+l+1-\lambda}{k[(k+1)+2(l+1)]}.
\label{eq:hydrogen-laguerre-recurrence}
\]
当且仅当 $\lambda$ 是满足 $n\ge l+1$ 的整数 $n$ 时,这个级数会终止,从而给出式~\eqref{eq:hydrogen-laguerre-equation} 的多项式解.于是可由 $\lambda$ 求出能量:
\[
 |E|=\frac{\mu Q^4}{2\hbar^2n^2}.
\]
记住 $E$ 为负数,就得到定理中给出的能级形式;条件 $n\ge l+1$ 与 $l\le n-1$ 相同.最后,把 $\rho$ 关于 $r$ 的公式和 $f$ 关于 $g$ 的公式代回,即得到定理所述的解形式.
\end{proof}

必须强调,定理~\ref{thm:hydrogen-bound-states} 中的函数并不张成整个 Hilbert 空间 $L^2(\mathbb R^3)$.毕竟这些函数都是 $\widehat H$ 的负特征值特征向量;如果它们张成整个空间,那么能量的期望值总会是负的.但很容易构造出属于 $\widehat H$ 定义域且满足 $\langle\psi,\widehat H\psi\rangle>0$ 的函数:只需取一个平均位置远离原点、平均动量很大的 Gaussian 波包.此时 $\langle\psi,V\psi\rangle$ 接近零,而 $\langle\psi,P^2\psi\rangle$ 很大且为正.

尽管如此,可以证明定理~\ref{thm:hydrogen-bound-states} 中的函数张成 $L^2(\mathbb R^3)$ 的负能量子空间.还可以分析 $\widehat H$ 的正谱部分,不过零以上的谱是纯连续谱,表示氢原子发生电离,即电子已经逃离质子.

\begin{theorem}[氢原子的负能量子空间]\label{thm:hydrogen-negative-spectrum}
当 $n$ 遍历所有正整数,$l$ 遍历 $0,1,\ldots,n-1$,而 $g$ 遍历所有次数为 $l$ 的齐次调和多项式时,定理~\ref{thm:hydrogen-bound-states} 中的特征向量张成 $L^2(\mathbb R^3)$ 的负能量子空间,即由谱定理给出的投影值测度 $\mu_{\widehat H}$ 的投影
\[
 \mu_{\widehat H}((-\infty,0))
\]
的值域.
\end{theorem}

\begin{proof}
证明需要用到超出第 9、10 章所发展工具的谱理论结果,本书无法完整重现.具体而言,我们使用文献 [27] 的定理 V.5.7,它说明 $\widehat H$ 的负能量部分是离散的,由有限重特征值组成,并且这些特征值只能在 $0$ 处聚积.

下面简要说明为什么上述结果成立.若 $A,B$ 是无界自伴算子,称 $B$ 是 $A$ 的相对紧扰动,如果对 $B$ 的预解集中的每个 $\lambda$,算子 $A(B-\lambda I)^{-1}$ 都是紧算子.根据文献 [27] 的引理 V.5.8,氢原子的势能算子是动能算子的相对紧扰动.这强化了我们在定理~\ref{thm:kato-rellich} 的证明中得到的结论,即势能算子关于动能算子相对有界,且相对界小于 $1$.相对紧性的证明依赖于氢原子的势在无穷远处趋于零这一事实.

另一方面,如果 $\lambda$ 是无界自伴算子 $A$ 的本质谱,就说 $\lambda$ 属于 $A$ 的本质谱;这意味着 $\lambda$ 要么是 $\sigma(A)$ 中的非孤立点,要么是 $A$ 的无限重特征值.根据文献 [27] 的定理 IV.5.35,自伴算子的相对紧扰动不改变本质谱.因此 $\widehat H$ 的本质谱等于动能算子的本质谱;由于动能算子非负,这个本质谱显然包含于 $[0,\infty)$.于是 $\widehat H$ 负能量部分的任何谱点都必须是 $\sigma(\widehat H)$ 中的孤立点,并且是有限重特征值.

由此,$\widehat H$ 在零以下没有连续谱,我们只需寻找平方可积特征向量.此外,$\widehat H$ 的每个 $E<0$ 特征子空间都是有限维的,因此它分解为不可约的 $SO(3)$ 不变子空间的直和.根据命题~\ref{prop:radial-representation-classification},这类子空间总是定义~\ref{def:radial-representation-subspace} 中某个 $V_{l,f}$ 的形式.因此,我们寻找形如 $\psi(x)=p(x)f(|x|)$ 的函数,使 $\widehat H\psi=E\psi$,其中 $E<0$.

如果形如 $p(x)f(|x|)$ 的函数是 Hamilton 算子的特征函数,则 $f$ 必须满足微分方程~\eqref{eq:hydrogen-radial-equation}.根据线性常微分方程的初等结果,对任意 $E$,这个方程恰有两个线性无关的解.两个解都可以设成式~\eqref{eq:hydrogen-radial-ansatz} 的形式,引入变量 $\rho$,再对 $g(\rho)$ 作幂级数展开来构造(见习题 9).其中一个 $g(\rho)$ 的幂级数以 $\rho^{-(2l+1)}$ 开头,这时 $\psi(x)$ 在原点附近像 $1/|x|^{l+1}$ 一样发散;这样的函数不属于 Hamilton 算子的定义域(见第 9 章习题 14).另一个解的级数从 $\rho^0$ 项开始,可以用式~\eqref{eq:hydrogen-radial-ansatz},把变量从 $r$ 换为 $\rho$,再用递推关系~\eqref{eq:hydrogen-laguerre-recurrence} 定义系数.如果所得级数不终止,那么不难看出,当 $k$ 很大时,其项的行为类似于 $e^\rho$ 的级数.由于 $f=e^{-\rho/2}g(\rho)$,这个函数在无穷远处将像 $e^{\rho/2}$ 一样增长,从而 $\psi$ 不属于 $L^2(\mathbb R^3)$.因此,要得到平方可积解,$g(\rho)$ 的级数必须终止;此时 $\psi$ 正是定理~\ref{thm:hydrogen-bound-states} 中的函数.
\end{proof}

\begin{corollary}[氢原子能级的简并度]\label{cor:hydrogen-energy-multiplicity}
定理~\ref{thm:hydrogen-bound-states} 中给出的每个能级 $E_n$ 的重数为 $n^2$.
\end{corollary}

\begin{proof}
根据定理~\ref{thm:hydrogen-negative-spectrum},定理~\ref{thm:hydrogen-bound-states} 中的特征向量构成 $\widehat H$ 特征值为 $E_n$ 的全部特征向量.因此,特征值 $E_n$ 对应的线性无关特征向量数目,就是 $l=0,1,\ldots,n-1$ 时球谐空间 $V_l$ 的维数之和.由定理~\ref{thm:spherical-harmonics-decomposition},这个数为
\[
 \sum_{l=0}^{n-1}(2l+1)=n^2,
\]
结论得证.
\end{proof}
}

\ChapterEighteenRadialRemainder
\ChapterEighteenHydrogen
\ChapterEighteenClassical
\ChapterEighteenQuantumRungeLenz
\ChapterEighteenSOFourSpin
\ChapterEighteenComputations
\ChapterEighteenLateCalculations
\ChapterEighteenExercises
